% !PRINTED 390
\weberpage{390}

To prove this, it suffices to show that by the permutations
of this group any root, say [1\ 3], can pass into
every other (with the exception of [1\ 2]). Now
[1,\,3] % [sic] comma printed here, unlike the bracket pairs elsewhere on this page which are space-separated without a comma
passes, by permutations of the group $S$ by which
1, 2 remain unchanged, into [1\ 4], \dots, [1\ 8]; likewise [2\ 3] into
[2\ 4], \dots, [2\ 8]. It therefore remains to show that [1\ 3]
can also be carried into [2\ 3] and into [4\ 5]. But this is shown
by inspection of the three following complete systems of seven:

\[
\begin{aligned}
&[1\ 2]\ [1\ 3]\ [1\ 4]\ [1\ 5]\ [1\ 6]\ [1\ 7]\ [1\ 8]\\
&[1\ 2]\ [2\ 3]\ [2\ 4]\ [2\ 5]\ [2\ 6]\ [2\ 7]\ [2\ 8]\\
&[1\ 2]\ [4\ 5]\ [3\ 4]\ [3\ 5]\ [1\ 6]\ [1\ 7]\ [1\ 8].
\end{aligned}
\]

According to this the group $P$ is doubly transitive. But it
cannot be more than doubly transitive. For if one leaves
two roots unchanged, a root syzygetic with these two
cannot be carried into an azygetic one.
By the adjunction of two roots the double tangent equation
becomes reducible. A factor of degree $10$ separates off,
whose roots, together with the two given ones, form a \gesperrt{Steiner}
complex (the function $X$ in Theorem 1., \S.\,101).

Among the divisors of the group $P$ the
group $S$ of index 36 is especially remarkable. This group is still
transitive, because by it [1\ 2] can be carried into any arbitrary root $[\varkappa\ \lambda]$.
But it is only simply transitive,
because when [1\ 2] is held fixed the root [1\ 3] can no longer
pass into [4\ 5].

By the adjunction of \gesperrt{one} root of an equation of degree $36$,
the group of the double tangent equation is therefore reduced to the
group of a general equation of degree $8$.

If one sets, e.\,g.,

\[
v_1 = [1\ 2]\ [1\ 3]\ [1\ 4]\ [1\ 5]\ [1\ 6]\ [1\ 7]\ [1\ 8],
\]

and defines $v_2, v_3, \dots, v_8$ correspondingly, then every symmetric
function of these eight quantities, e.\,g.

\[
u = v_1 + v_2 + v_3 + v_4 + v_5 + v_6 + v_7 + v_8,
\]

is a root of an equation of degree $36$. If one adjoins to the
problem the quantity $u$, the $v_1, v_2, \dots, v_8$ become the roots
of an equation of degree $8$, which has no affect.
